% One run of the harness. The four variants differ only in which of the three
% shaded states is entered.
\begin{tikzpicture}[font=\sffamily\scriptsize,
  opt/.style={state, fill=hwcol, text width=30mm, minimum height=7mm,
              font=\sffamily\scriptsize}]

\node[initial] (i) at (0,0) {};
\node[rtstate] (fill) at (1.8,0) {fill 0xA5};
\node[rtstate] (arm) at (4.6,0) {arm watermark};
\node[rtstate] (wint) at (7.4,0) {wait for INT1};
\node[rtstate] (start) at (10.2,0) {start transfer};
\node[rtstate] (wc) at (13.0,0) {wait complete};

\draw[trans] (i) -- (fill);
\draw[trans] (fill) -- (arm);
\draw[trans] (arm) -- (wint);
\draw[trans] (wint) -- node[above, font=\tiny]{edge} (start);
\draw[trans] (start) -- (wc);

\node[opt] (o1) at (12.6,-2.2) {invalidate by address};
\node[opt] (o2) at (12.6,-3.8) {nothing: the region is non-cacheable};
\node[opt] (o3) at (12.6,-5.4) {nothing: the buffer is not cached};

\draw[trans] (wc.south) -- node[right, font=\tiny, xshift=1mm]{maintain} (o1.north);
\draw[trans] (o1.west) -- ++(-0.5,0) |- (o2.west);
\draw[trans] (o2.west) -- ++(-0.5,0) |- (o3.west);

\node[rtstate] (count) at (7.6,-6.8) {count survivors};
\node[rtstate] (acc) at (4.4,-6.8) {accumulate};
\node[rtstate] (rep) at (1.4,-6.8) {report rate};
\node[final] (f) at (1.4,-8.6) {};

\draw[trans] (wc.east) -- ++(0.9,0)
  node[right, font=\tiny, align=left, xshift=-1mm, yshift=-3mm]{naive}
  |- (count.east);
\draw[trans] (o3.south) |- (count.east);
\draw[trans] (count) -- (acc);
\draw[trans] (acc) -- node[above, font=\tiny, fill=white, inner sep=1pt]{1000 done} (rep);
\draw[trans] (rep) -- (f);
\draw[trans] (acc.north) -- ++(0,0.9)
  node[above, font=\tiny, xshift=8mm]{not yet 1000} -| (fill.south);

\node[umlnote, text width=56mm, anchor=north west] at (0,-9.6)
  {Everything outside the shaded band is identical in all four variants. That
   is what makes the four failure counts comparable without argument, and it is
   the reason the harness is deliberately dull.};
\node[umlnote, text width=56mm, anchor=north west] at (7.6,-9.6)
  {A surviving sentinel byte means the processor read a cache line that memory
   had already moved past. The naive variant is required to produce some, or
   the experiment did not reproduce the condition.};
\end{tikzpicture}
